\documentclass[10pt,
% a4paper,
%twocolumn,
fleqn,
%landscape, 
% papersize,
dvipdfmx,
uplatex
]{jsarticle}



\def\maru#1{\textcircled{\scriptsize#1}}%丸囲み番号

% \RequirePackage[2020/09/30]{platexrelease}

%太字設定
\usepackage[deluxe]{otf}

\usepackage{emathEy}

\usepackage[g]{esvect}

%大きな文字
\usepackage{fix-cm}

%定理環境
\usepackage{emathThm}
%\theoremstyle{boxed}
\theorembodyfont{\normalfont}
\newtheorem{Question}{問題}[subsection]
\newtheorem{Q}{}[subsection]
\newtheorem{question}[Question]{}
\newtheorem{quuestion}{}[subsection]

%セクション，大問番号のデザイン
\renewcommand{\labelenumi}{(\arabic{enumi})}
\renewcommand{\theenumii}{\alph{enumii})}
\renewcommand{\thesection}{第\arabic{section}章}

%用紙サイズの詳細設定
\usepackage{bxpapersize}
\papersizesetup{size={80mm,45mm}}
\usepackage[top=0.7zw,bottom=0truemm,left=3truemm,right=133truemm]{geometry}
\usepackage[dvipdfmx]{graphicx}

%余白など
\usepackage{setspace} % 行間
\setlength{\mathindent}{1zw}
\setlength\parindent{0pt}


%色カラーに関する設定
\usepackage{color}
\definecolor{shiro}{rgb}{0.95703125,0.87109375,0.7421875}
\definecolor{kin}{rgb}{0.95703125,0.87109375,0.7421875}
\definecolor{orange}{rgb}{1,0.7,0.2}
\definecolor{bradorange}{rgb}{1,0.5,0}
\definecolor{pink}{rgb}{0.9176,0.5686,0.5960}
\definecolor{mizu}{rgb}{0.6156,0.8,0.9955}
\color{kin}
% \pagecolor{hukamido}

\usepackage{at}%図の配置
% \usepackage{wallpaper}

\begin{document}



\at(0cm,0cm){\includegraphics[width=8cm,bb=0 0 1920 1080]{./youtube/thumbnails/templates/smart_background/２次関数.jpeg}}
{\color{orange}\bf\boldmath\huge\underline{放物線の平行移動}}\vspace{0.3zw}

\large
\bf\boldmath 問．放物線$y=2x^2-8x+5$を

\huge
\vspace{-0.1zw}
\hspace{0.1zw}
どのように\vspace{-0.1zw}\\
\hfill 平行移動すると
\vspace{0.2zw}
\hspace{0.1zw}

\large
\hfill 放物線$y=2x^2+4x+7$に重なるか．
\at(7.0cm,0.2cm){\small\color{bradorange}$\overset{\text{２次関数}}{\text{基本}}$}


\newpage



\at(0cm,0cm){\includegraphics[width=8cm,bb=0 0 1920 1080]{./youtube/thumbnails/templates/smart_background/２次関数.jpeg}}
{\color{orange}\bf\boldmath\huge\underline{放物線の対称移動}}\vspace{0.3zw}

\large 
\bf\boldmath 問．放物線$y=3x^2-6x+4$を

\LARGE
\hspace{0.1zw}
$x$軸，$y$軸，原点関して\\
\hfill 対称移動した放物線
\hspace{0.1zw}

\large
\vspace{0.5zw}
\hfill の方程式をそれぞれ求めよ．
\at(7.0cm,0.2cm){\small\color{bradorange}$\overset{\text{２次関数}}{\text{基本}}$}


\newpage



\at(0cm,0cm){\includegraphics[width=8cm,bb=0 0 1920 1080]{./youtube/thumbnails/templates/smart_background/２次関数.jpeg}}
{\color{orange}\bf\boldmath\huge\underline{放物線の平行移動}}\vspace{0.3zw}

\large 
\bf\boldmath 問．放物線$y=3x^2-6x+4$を

\LARGE
\hspace{-0.1zw}
$x$軸方向に$-1,y$軸方向に$3$\\
\hfill だけ平行移動した放物線
\hspace{0.1zw}

\large
\vspace{0.5zw}
\hfill の方程式を求めよ．
\at(7.0cm,0.2cm){\small\color{bradorange}$\overset{\text{２次関数}}{\text{基本}}$}


\newpage



\at(0cm,0cm){\includegraphics[width=8cm,bb=0 0 1920 1080]{./youtube/thumbnails/templates/smart_background/２次関数.jpeg}}
{\color{orange}\bf\boldmath\huge\underline{$2$次関数のグラフ}}\vspace{0.3zw}

\Large 
\bf\boldmath 問．次の$2$次関数のグラフをかけ．

\huge
(1)  $y=x^2+4x+1$\\
(2)  $y=-2x^2+4x+3$\\

\at(7.0cm,0.2cm){\small\color{bradorange}$\overset{\text{２次関数}}{\text{基本}}$}


\newpage



\at(0cm,0cm){\includegraphics[width=8cm,bb=0 0 1920 1080]{./youtube/thumbnails/templates/smart_background/２次関数.jpeg}}
{\color{orange}\bf\boldmath\huge\underline{グラフの拡大}}\vspace{0.1zw}

\normalsize
\bf\boldmath 問．関数$y=2x^2$のグラフを$y$軸をもとにして，

\huge
\vspace{-0.0zw}
\hspace{0.1zw}
$x$軸方向に\\
\hfill $2$倍に拡大する
\hspace{0.1zw}

\normalsize
\vspace{0.5zw}
\hfill ことで得られるグラフの式を求めよ．
\at(7.0cm,0.2cm){\small\color{bradorange}$\overset{\text{２次関数}}{\text{基本}}$}


\newpage



\at(0cm,0cm){\includegraphics[width=8cm,bb=0 0 1920 1080]{./youtube/thumbnails/templates/smart_background/２次関数.jpeg}}
{\color{orange}\bf\boldmath\huge\underline{グラフの平行移動}}\vspace{0.3zw}

\large 
\bf\boldmath 問．関数$y=2x^2$のグラフを

\LARGE
\hspace{0.1zw}
$x$軸方向に$3,\;y$軸方向に$1$\\
\hfill だけ平行移動する
\hspace{0.1zw}

\large
\vspace{0.5zw}
\hfill ことで得られるグラフの式を求めよ．
\at(7.0cm,0.2cm){\small\color{bradorange}$\overset{\text{２次関数}}{\text{基本}}$}


\newpage



\at(0cm,0cm){\includegraphics[width=8cm,bb=0 0 1920 1080]{./youtube/thumbnails/templates/smart_background/２次関数.jpeg}}
{\color{orange}\bf\boldmath\LARGE\underline{$2$次方程式 $3$つの解法}}\vspace{0.3zw}

\normalsize
\bf\boldmath 問．次の$2$次方程式を解け．

\LARGE
\vspace{-0.2zw}
(1)  $2x^2-5x-3=0$\\
(2)  $x^2-3=0$\\
(3)  $x^2-4x+1=0$\\

\at(7.0cm,0.2cm){\small\color{bradorange}$\overset{\text{２次関数}}{\text{基本}}$}


\newpage



\at(0cm,0cm){\includegraphics[width=8cm,bb=0 0 1920 1080]{./youtube/thumbnails/templates/smart_background/２次関数.jpeg}}
{\color{orange}\bf\boldmath\huge\underline{解の公式の証明}}\vspace{0.1zw}

\Large 
\bf\boldmath 問．$2$次方程式の解の公式


\vspace{0.5zw}
\hspace{0.0zw}$x=$
\huge$\bunsuu{{-b+\sqrt {b^2-4ac}}}{{2a}}\vspace{0.3zw}$

\Large 
\hfill を証明せよ．
\at(7.0cm,0.2cm){\small\color{bradorange}$\overset{\text{２次関数}}{\text{基本}}$}


\newpage



\at(0cm,0cm){\includegraphics[width=8cm,bb=0 0 1920 1080]{./youtube/thumbnails/templates/smart_background/２次関数.jpeg}}
{\color{orange}\bf\boldmath\huge\underline{判別式$D$ 全パターン}}\vspace{-0.1zw}

\huge 
\bf\boldmath 問． 次の$2$次方程式の\\
\hfill 実数解の個数を求めよ．

\normalsize 
\vspace{0.4zw}
\hspace{-1zw}
$\begin{array}{ll}
(1)\hspace{0.3zw}  x^2-4x+1=0&\hspace{-0.5zw}
(2)\hspace{0.3zw}  2x^2-x+4=0\\
(3)\hspace{0.3zw}  4x^2+{12}x+9=0&\hspace{-0.5zw}
(4)\hspace{0.3zw}  2x^2-x-4=0\\
\end{array}$

\at(7.0cm,0.2cm){\small\color{bradorange}$\overset{\text{２次関数}}{\text{基本}}$}


\newpage



\at(0cm,0cm){\includegraphics[width=8cm,bb=0 0 1920 1080]{./youtube/thumbnails/templates/smart_background/２次関数.jpeg}}
{\color{orange}\bf\boldmath\huge\underline{$2$次不等式 全パターン}}\vspace{0.5zw}

\large
\bf\boldmath 問．次の$2$次不等式を解け．

\small
\vspace{0.8zw}
\hspace{-1zw}
$\begin{array}{ll}
(1)\hspace{0.3zw}  x^2-7x+6\leqq 0&\hspace{-0.5zw}
(2)\hspace{0.3zw}  x^2-2x-{15}>0\vspace{0.5zw}\\
(3)\hspace{0.3zw}  2x^2+x-6<0&\hspace{-0.5zw}
(4)\hspace{0.3zw}  -x^2+2x+4\leqq 0\vspace{0.5zw}\\
(5)\hspace{0.3zw}  x^2-{12}x+{36}>0&\hspace{-0.5zw}
(6)\hspace{0.3zw}  9x^2+{24}x+{16}\leqq 0\vspace{0.5zw}\\
(7)\hspace{0.3zw}  x^2-x+3>0&\hspace{-0.5zw}
(8)\hspace{0.3zw}  -2x^2-6x-5>0\vspace{0.5zw}\\
\end{array}$

\at(7.0cm,0.2cm){\small\color{bradorange}$\overset{\text{２次関数}}{\text{基本}}$}


\end{document}

